%% ----------------------------------------------------------------
%% Codes and Standards
%% ----------------------------------------------------------------
\phantomsection
\addtotoc{Codes and Standards}
\chapter*{Codes and Standards}
\thispagestyle{plain}

\noindent\textbf{Project Title:} \textit{Asynchronous Sensor Fusion for ISAC Nodes Integrating Wi-Fi Sensing and mmWave Radar}

\vspace{0.8cm}

\noindent\textbf{1. Wireless Communication \& Sensing Standards}

\medskip
\noindent\textbf{IEEE 802.11mc -- Fine Time Measurement (FTM)}\\
IEEE standard that defines how a round trip time (RTT) measurement is made using MAC address information from the Wi-Fi responder.

\medskip
\noindent\textbf{IEEE 802.11 b/g/n -- 2.4~GHz Wi-Fi Standard}\\
Standard that provides the basic communication structure over which all other standards operate; this particular version of the standard allows for high-throughput (HT) operation at a 40~MHz bandwidth, allowing for higher-resolution timing.

\vspace{0.6cm}
\noindent\textbf{2. Radar \& Frequency Band Standards}

\medskip
\noindent\textbf{ITU-R/WPC (India) Regulations for ISM Bands}\\
Regulation defining the use of FMCW radar systems that operate at 24~GHz frequency modulation continuous wave in order to meet local RF emission regulations and bandwidth requirements as defined under the Indian government's industrial scientific medical (ISM) bands.

\vspace{0.6cm}
\noindent\textbf{3. Embedded Systems \& Interfacing Standards}

\medskip
\noindent\textbf{UART (Universal Asynchronous Receiver-Transmitter) Protocol}\\
Protocol that defines how an embedded system uses serial communication to communicate data between two or more embedded systems at a speed of 256~kbaud. The HLK-LD2450 radar module interfaces with a microcontroller via this method.

\medskip
\noindent\textbf{FreeRTOS Task Scheduling Standards}\\
Set of guidelines that define the operational characteristics of a real-time operating system (RTOS), including task scheduling and isolation, interrupt handling and inter-process communication (IPC). This set of guidelines applies to both cores of the dual-core processor.

\vspace{0.6cm}
\noindent\textbf{4. Machine Learning \& Software Standards}

\medskip
\noindent\textbf{TensorFlow Lite Micro Guidelines}\\
Provides a standard definition for deploying Float32 multi-layer perceptron (MLP) machine learning models onto resource constrained edge computing devices such that they can run completely contained within the 32~KB static RAM without requiring additional memory outside the SRAM.
